\chapter[Свойства интеграла с переменным верхним пределом (непрерывность, дифференцируемость). Формула Ньютона"--~Лейбница.]{Свойства интеграла с переменным верхним пределом (непрерывность, дифференцируемость). Формула Ньютона"--~Лейбница.}
\section{Определённый интеграл Римана}

\subsection{Разбиения, суммы Дарбу и интегральные суммы}

Напомним, что разбиением отрезка $[a, b]$ называется конечный набор точек $\mathrm{T} = \{x_i\}_{i=0}^I$, таких, что $a = x_0 < x_1 < \ldots < x_I = b$. Отрезки $[x_{i-1}, x_i]$ называются отрезками разбиения $\mathrm{T}$.

\begin{defn}
Пусть на $[a, b]$ определена функция $f(x)$ и задано разбиение $\mathrm{T} = \{x_i\}_{i=0}^I$ отрезка $[a, b]$. Определим
$$
m_i = \inf_{x \in [x_{i-1}, x_i]} f(x), \qquad M_i = \sup_{x \in [x_{i-1}, x_i]} f(x),
$$
$$
s(f; \mathrm{T}) = \sum_{i=1,\dots,I} (x_i - x_{i-1})m_i,
$$
$$
S(f; \mathrm{T}) = \sum_{i=1,\dots,I} (x_i - x_{i-1})M_i.
$$
Сумма $s(f; \mathrm{T})$ называется \textit{нижней суммой Дарбу}, а $S(f; \mathrm{T})$ "--- \textit{верхней суммой Дарбу} для функции $f$ и разбиения $\mathrm{T}$.
\end{defn}

\begin{defn}
\textit{Мелкостью разбиения} $\mathrm{T} = \{x_i\}_{i=0}^I$ называется число
$$
\ell(\mathrm{T}) = \max_{i=1,\dots,I} (x_i - x_{i-1}).
$$
\end{defn}

\begin{defn}
Число $J$ называется (\textit{определенным}) \textit{интегралом Римана} функции $f$ на $[a, b]$ и обозначается $J = \int\limits_a^b f(x)\,dx$, если
$$
\lim_{\ell(\mathrm{T}) \to 0} s(f; \mathrm{T}) = \lim_{\ell(\mathrm{T}) \to 0} S(f; \mathrm{T}) = J, \text{ т.\,е.}
$$
$$
\forall\, \varepsilon > 0 \;\; \exists\, \delta > 0 \;\; \forall\, \mathrm{T} \colon \;\; \ell(\mathrm{T}) \le \delta \hookrightarrow |s(f; \mathrm{T}) - J| \le \varepsilon \;\; \text{и} \;\; |S(f; \mathrm{T}) - J| \le \varepsilon.
$$
Функция $f$ называется \textit{интегрируемой по Риману} на $[a, b]$, если существует интеграл Римана функции $f$ на $[a, b]$.
\end{defn}

\begin{defn}
Пусть задано разбиение $\mathrm{T} = \{x_i\}_{i=0}^I$ отрезка $[a, b]$. \textit{Выборкой}, соответствующей разбиению $\mathrm{T}$, называется набор точек $\xi_{\mathrm{T}} = \{\xi_i\}_{i=1}^I$ таких, что $\xi_i \in [x_{i-1}, x_i]$. \textit{Интегральной суммой (Римана)} для функции $f$, разбиения $\mathrm{T}$ и выборки $\xi_{\mathrm{T}}$ называется
$$
\sigma(f; \mathrm{T}; \xi_{\mathrm{T}}) = \sum_{i=1}^I (x_i - x_{i-1})f(\xi_i).
$$
\end{defn}


\section{Интеграл с переменным верхним пределом и формула Ньютона "--- Лейбница}

\subsection{Свойства интеграла как функции верхнего предела}

\begin{thm}[непрерывность интеграла как функции верхнего предела] \label{th:int_cont}
Пусть на отрезке $[a, b]$ задана интегрируемая по Риману функция $f(x)$. Тогда функция $F(x) = \int\limits_a^x f(t)\,dt$ непрерывна на $[a, b]$.
\end{thm}

\begin{proof}
В силу необходимого условия интегрируемости функция $f(x)$ ограничена на $[a, b]$, т.\,е.
$$
\exists\, C \in \bbR \colon \;\; \forall\, x \in [a, b] \hookrightarrow |f(x)| \le C.
$$
Пусть $x_1, x_2 \in [a, b]$. В силу свойства аддитивности интеграла относительно отрезков интегрирования $F(x_2) - F(x_1) = \int\limits_{x_1}^{x_2} f(t)\,dt$. По теореме об интегрировании неравенств
$$
|F(x_2) - F(x_1)| \le \left| \int\limits_{x_1}^{x_2} C\,dt \right| = C|x_2 - x_1|.
$$
Следовательно,
$$
\forall\, x_0 \in [a, b] \;\; \forall\, \varepsilon > 0 \;\; \exists\, \delta = \frac{\varepsilon}{C} \;\; \forall\, x \in [a, b] \colon \;\; |x - x_0| < \delta \hookrightarrow |F(x) - F(x_0)| < \varepsilon,
$$
т.\,е. функция $F(x)$ непрерывна на $[a, b]$.
\end{proof}

\begin{defn}
Функция $F(x)$ называется \textit{первообразной функцией} $f(x)$ на $[a, b]$, если $\forall\, x \in (a, b) \hookrightarrow F'(x) = f(x)$, а на концах отрезка $[a, b]$ значения функции $f$ равны односторонним производным функции $F$:
$$
f(a) = F'_+(a) = \lim_{x \to a+0} \frac{F(x) - F(a)}{x - a}, \qquad f(b) = F'_-(b) = \lim_{x \to b-0} \frac{F(x) - F(b)}{x - b}.
$$
\end{defn}

\begin{thm} \label{th:int_antideriv}
Если функция $f$ непрерывна на $[a, b]$, то функция $F(x) = \int\limits_a^x f(t)\,dt$ является первообразной функции $f(x)$ на $[a, b]$.
\end{thm}

\begin{proof}
Пусть $x, x_0 \in [a, b]$, $x \ne x_0$. Тогда
$$
F(x) - F(x_0) = \int\limits_{x_0}^x f(t)\,dt = (x - x_0)f(x_0) + \int\limits_{x_0}^x (f(t) - f(x_0))\,dt.
$$
Следовательно,
$$
\left| \frac{F(x) - F(x_0)}{x - x_0} - f(x_0) \right| = \left| \frac{1}{x - x_0} \int\limits_{x_0}^x (f(t) - f(x_0))\,dt \right| \le
$$
$$
\le \frac{1}{|x - x_0|} \left| \int\limits_{x_0}^x |f(t) - f(x_0)|\,dt \right|.
$$
В силу непрерывности функции $f$ на $[a, b]$
$$
\forall\, \varepsilon > 0 \;\; \exists\, \delta > 0 \;\; \forall\, t \in [a, b] \colon \;\; |t - x_0| \le \delta \hookrightarrow |f(t) - f(x_0)| \le \varepsilon,
$$
поэтому
$$
\forall\, \varepsilon > 0 \;\; \exists\, \delta > 0 \;\; \forall\, x \in [a, b] \colon \;\; |x - x_0| \le \delta \hookrightarrow \left| \int\limits_{x_0}^x |f(t) - f(x_0)|\,dt \right| \le |x - x_0|\varepsilon.
$$
Следовательно,
$$
\forall\, \varepsilon > 0 \;\; \exists\, \delta > 0 \;\; \forall\, x \in [a, b] \colon \;\; |x - x_0| \le \delta \hookrightarrow \left| \frac{F(x) - F(x_0)}{x - x_0} - f(x_0) \right| \le \varepsilon,
$$
то есть $\forall\, x_0 \in [a, b] \hookrightarrow \lim\limits_{x \to x_0} \frac{F(x) - F(x_0)}{x - x_0} = f(x_0)$, где при $x_0 = a$ имеется в виду предел справа, а при $x_0 = b$ "--- предел слева. Это означает, что $F'_+(a) = f(a)$, $F'_-(b) = f(b)$, $\forall\, x_0 \in (a, b) \hookrightarrow F'(x_0) = f(x_0)$. Таким образом, функция $F$ является первообразной функции $f$ на $[a, b]$.
\end{proof}


\subsection{Основная формула интегрального исчисления}

\begin{cons} \label{cons:general_antideriv}
Любая первообразная непрерывной на $[a, b]$ функции $f$ имеет вид
$$
F(x) = \int\limits_a^x f(t)\,dt + C,
$$
где $C \in \bbR$ "--- произвольная константа.
\end{cons}

\begin{cons}[формула Ньютона"---~Лейбница] \label{cons:newton_leibniz}
Если $F$ "--- первообразная непрерывной на $[a, b]$ функции $f$, то
$$
\int\limits_a^b f(x)\,dx = F(x) \bigg|_a^b \stackrel{\text{по опред.}}{=} F(b) - F(a).
$$
\end{cons}

\begin{proof}
Воспользуемся следствием \ref{cons:general_antideriv} и заметим, что
$$
F(a) = \int\limits_a^a f(t)\,dt + C = C, \qquad F(b) = \int\limits_a^b f(t)\,dt + C = \int\limits_a^b f(t)\,dt + F(a).
$$
Следовательно, $\int\limits_a^b f(x)\,dx = F(b) - F(a)$.
\end{proof}